% GENERATED FILE. Edit canonical lessons, then run npm run generate:books.
% canonical-source-hash: fnv1a64:911982cafac87a2a
% canonical-lessons: HI-C312-rediyo, HI-C312-samacar, HI-C312-prasaran, HI-C312-surkhi, HI-C312-chapna

\chapter{The Morning News}
\label{ch:hi-the-morning-news}
\hlchaptermodality{312}

\begin{chapteropening}
\textbf{By the end of this chapter:} \emph{I can say what the news reported.}

Complete the last lesson of chapter 312: I can say what the news reported.
\end{chapteropening}

\section[\texorpdfstring{\hi{रेडियो} (reḍiyo)}{रेडियो (reḍiyo)}]{\hi{रेडियो} (reḍiyo) — a radio}

\label{lesson:HI-C312-rediyo}



\begin{quote}
Before the new one: say the Hindi for to collect, then the Hindi for spirit.
\end{quote}

\subsection*{You'll want to know: \hi{रेडियो}}
\textbf{\hi{रेडियो}} — \emph{reḍiyo} — \textquotedblleft{}a radio\textquotedblright{}.

\begin{quote}
\hi{दादा} \hi{जी} \hi{हर} \hi{सुबह} \hi{रेडियो} \hi{सुनते} \hi{हैं।}
\end{quote}

\begin{quote}
\emph{dādā jī har subah reḍiyo sunte haiṁ.}
\end{quote}

\begin{quote}
Grandfather listens to the radio every morning.
\end{quote}

Masculine.

\begin{cousinweb}
\textbf{\hi{रेडियो}} (\emph{reḍiyo}) is borrowed from English \emph{radio}, taken into Hindi as a masculine noun: \textbf{\hi{रेडियो} \hi{पर}}, on the radio.
\end{cousinweb}

\begin{grammarlens}[title={What you've built}]
One more word for telling what happened.
\end{grammarlens}

\subsection*{Your turn}
\begin{itemize}
\raggedright
  \item \emph{Say it:} \emph{reḍiyo}
  \item \emph{Say it:} \emph{reḍiyo}, once more
  \item \emph{Say it:} \emph{dādā jī har subah reḍiyo sunte haiṁ.} — Grandfather listens to the radio every morning.
  \item \emph{Recall:} say the Hindi for to collect, then the Hindi for spirit, then say \emph{reḍiyo} again
\end{itemize}

\begin{tcolorbox}[breakable,colback=teal!4,colframe=teal!35!black,title={Before you move on}]
What does \hi{रेडियो} mean? (\textquotedblleft{}A radio\textquotedblright{}.) Say it once more.
\end{tcolorbox}
\section[\texorpdfstring{\hi{समाचार} (samācār)}{समाचार (samācār)}]{\hi{समाचार} (samācār) — the news (formal); a news report}

\label{lesson:HI-C312-samacar}



\begin{quote}
Before the new one: say the Hindi for spirit, then the Hindi for a radio.
\end{quote}

\subsection*{You'll want to know: \hi{समाचार}}
\textbf{\hi{समाचार}} — \emph{samācār} — \textquotedblleft{}the news (formal); a news report\textquotedblright{}.

\begin{quote}
\hi{रेडियो} \hi{पर} \hi{सुबह} \hi{आठ} \hi{बजे} \hi{समाचार} \hi{आते} \hi{हैं।}
\end{quote}

\begin{quote}
\emph{reḍiyo par subah āṭh baje samācār āte haiṁ.}
\end{quote}

\begin{quote}
The news is on the radio at eight in the morning.
\end{quote}

Masculine.

\begin{cousinweb}
\textbf{\hi{समाचार}} (\emph{samācār}) is Sanskrit \emph{samācāra} (Monier-Williams), first conduct and practice and then news of what someone has done. It is the word of the broadcast; \textbf{\hi{ख़बर}} is the word of conversation.
\end{cousinweb}

\begin{grammarlens}[title={What you've built}]
One more word for telling what happened.
\end{grammarlens}

\subsection*{Your turn}
\begin{itemize}
\raggedright
  \item \emph{Say it:} \emph{samācār}
  \item \emph{Say it:} \emph{samācār}, once more
  \item \emph{Say it:} \emph{reḍiyo par subah āṭh baje samācār āte haiṁ.} — The news is on the radio at eight in the morning.
  \item \emph{Recall:} say the Hindi for spirit, then the Hindi for a radio, then say \emph{samācār} again
\end{itemize}

\begin{tcolorbox}[breakable,colback=teal!4,colframe=teal!35!black,title={Before you move on}]
What does \hi{समाचार} mean? (\textquotedblleft{}The news (formal); a news report\textquotedblright{}.) Say it once more.
\end{tcolorbox}
\section[\texorpdfstring{\hi{प्रसारण} (prasāraṇ)}{प्रसारण (prasāraṇ)}]{\hi{प्रसारण} (prasāraṇ) — a broadcast}

\label{lesson:HI-C312-prasaran}



\begin{quote}
Before the new one: say the Hindi for a radio, then the Hindi for the news (formal).
\end{quote}

\subsection*{You'll want to know: \hi{प्रसारण}}
\textbf{\hi{प्रसारण}} — \emph{prasāraṇ} — \textquotedblleft{}a broadcast\textquotedblright{}.

\begin{quote}
\hi{कल} \hi{रात} \hi{तूफ़ान} \hi{की} \hi{वजह} \hi{से} \hi{प्रसारण} \hi{बंद} \hi{रहा।}
\end{quote}

\begin{quote}
\emph{kal rāt tūfān kī vajah se prasāraṇ band rahā.}
\end{quote}

\begin{quote}
Last night the broadcast was off because of the storm.
\end{quote}

Masculine.

\begin{cousinweb}
\textbf{\hi{प्रसारण}} (\emph{prasāraṇ}) is Sanskrit \emph{prasāraṇa}, \textquotedblleft{}a spreading out, an extending\textquotedblright{} (Monier-Williams): news spread out over the air. The sense of broadcasting is modern.
\end{cousinweb}

\begin{grammarlens}[title={What you've built}]
One more word for telling what happened.
\end{grammarlens}

\subsection*{Your turn}
\begin{itemize}
\raggedright
  \item \emph{Say it:} \emph{prasāraṇ}
  \item \emph{Say it:} \emph{prasāraṇ}, once more
  \item \emph{Say it:} \emph{kal rāt tūfān kī vajah se prasāraṇ band rahā.} — Last night the broadcast was off because of the storm.
  \item \emph{Recall:} say the Hindi for a radio, then the Hindi for the news (formal), then say \emph{prasāraṇ} again
\end{itemize}

\begin{tcolorbox}[breakable,colback=teal!4,colframe=teal!35!black,title={Before you move on}]
What does \hi{प्रसारण} mean? (\textquotedblleft{}A broadcast\textquotedblright{}.) Say it once more.
\end{tcolorbox}
\section[\texorpdfstring{\hi{सुर्ख़ी} (surkhī)}{सुर्ख़ी (surkhī)}]{\hi{सुर्ख़ी} (surkhī) — a headline}

\label{lesson:HI-C312-surkhi}



\begin{quote}
Before the new one: say the Hindi for the news (formal), then the Hindi for a broadcast.
\end{quote}

\subsection*{You'll want to know: \hi{सुर्ख़ी}}
\textbf{\hi{सुर्ख़ी}} — \emph{surkhī} — \textquotedblleft{}a headline\textquotedblright{}.

\begin{quote}
\hi{आज} \hi{अख़बार} \hi{की} \hi{सबसे} \hi{बड़ी} \hi{सुर्ख़ी} \hi{बाढ़} \hi{पर} \hi{है।}
\end{quote}

\begin{quote}
\emph{āj akhbār kī sabse baṛī surkhī bāṛh par hai.}
\end{quote}

\begin{quote}
Today the newspaper's biggest headline is about the flood.
\end{quote}

Feminine.

\begin{cousinweb}
\textbf{\hi{सुर्ख़ी}} (\emph{surkhī}) is usually traced to Persian \emph{surkhī}, \textquotedblleft{}redness\textquotedblright{}, from \emph{surkh}, red; it is said that headings were once set in red. \textbf{\hi{सुर्ख़ियों} \hi{में} \hi{होना}} is to be in the headlines.
\end{cousinweb}

\begin{grammarlens}[title={What you've built}]
One more word for telling what happened.
\end{grammarlens}

\subsection*{Your turn}
\begin{itemize}
\raggedright
  \item \emph{Say it:} \emph{surkhī}
  \item \emph{Say it:} \emph{surkhī}, once more
  \item \emph{Say it:} \emph{āj akhbār kī sabse baṛī surkhī bāṛh par hai.} — Today the newspaper's biggest headline is about the flood.
  \item \emph{Recall:} say the Hindi for the news (formal), then the Hindi for a broadcast, then say \emph{surkhī} again
\end{itemize}

\begin{tcolorbox}[breakable,colback=teal!4,colframe=teal!35!black,title={Before you move on}]
What does \hi{सुर्ख़ी} mean? (\textquotedblleft{}A headline\textquotedblright{}.) Say it once more.
\end{tcolorbox}
\section[\texorpdfstring{\hi{छपना} (chapnā)}{छपना (chapnā)}]{\hi{छपना} (chapnā) — to be printed; to appear in print}

\label{lesson:HI-C312-chapna}



\begin{quote}
Before the new one: say the Hindi for a broadcast, then the Hindi for a headline.
\end{quote}

\subsection*{You'll want to know: \hi{छपना}}
\textbf{\hi{छपना}} — \emph{chapnā} — \textquotedblleft{}to be printed; to appear in print\textquotedblright{}.

\begin{quote}
\hi{हादसे} \hi{की} \hi{ख़बर} \hi{सब} \hi{अख़बारों} \hi{में} \hi{छपी।}
\end{quote}

\begin{quote}
\emph{hādse kī khabar sab akhbāroṁ meṁ chapī.}
\end{quote}

\begin{quote}
The news of the accident was printed in every newspaper.
\end{quote}

Past (it …): \hi{वह} \hi{ख़बर} \hi{छपी}, \emph{vah khabar chapī}.

\begin{cousinweb}
\textbf{\hi{छपना}} (\emph{chapnā}) is the intransitive partner of \textbf{\hi{छापना}}, to print, which CDIAL 4994 sets under \emph{*chapp}, a pressing: what is printed \textbf{\hi{छपता} \hi{है}}.
\end{cousinweb}

\begin{grammarlens}[title={What you've built}]
One more word for telling what happened.
\end{grammarlens}

\subsection*{Your turn}
\begin{itemize}
\raggedright
  \item \emph{Say it:} \emph{chapnā}
  \item \emph{Say it:} \emph{chapnā}, once more
  \item \emph{Say it:} \emph{hādse kī khabar sab akhbāroṁ meṁ chapī.} — The news of the accident was printed in every newspaper.
  \item \emph{Recall:} say the Hindi for a broadcast, then the Hindi for a headline, then say \emph{chapnā} again
\end{itemize}

\begin{tcolorbox}[breakable,colback=teal!4,colframe=teal!35!black,title={Before you move on}]
What does \hi{छपना} mean? (\textquotedblleft{}To be printed; to appear in print\textquotedblright{}.) Say it once more.
\end{tcolorbox}
